The final integrated QFlow core achieved timing closure at a 10~ns target after a staged optimization sequence in which the dominant bottleneck first resided in SKAG, then shifted to FDPE after SKAG repair, and was finally resolved by FDPE re-pipelining. The resulting tc5 design achieved approximately 101.885~MHz at post-synthesis, while OOC post-route implementation preserved positive setup and hold slack with an approximate frequency of 100.392~MHz.

For routing evaluation, the strongest non-trivial classical baseline was the key-aware shortest-path heuristic. A focused selection-policy study showed that an earlier collapse of the software PMO-GA reference to the key-aware heuristic was largely caused by latency-first final-answer selection. When the final candidate was instead chosen using weighted Tchebycheff scoring, PMO-GA became distinct from the key-aware baseline on a meaningful subset of the 6 evaluated cases. The recommended software PMO-GA variant is therefore \textbf{ga\_tcheby\_pick\_default}.

The dominant PMO-GA advantage in the current study is improved balance and load spreading rather than universal wins on every metric. The strongest balance gain appeared in \textbf{depletion\_plus\_stricter\_fidelity / irregular12}, with a load-balance change of -0.145170 relative to the key-aware baseline. The strongest fidelity-oriented case appeared in \textbf{current\_like / mesh16}, where fidelity improved by 0.000444. Accordingly, the honest claim is that QFlow supports \textbf{multi-objective tradeoff-aware routing with 100~MHz-class validated hardware closure}, rather than universal dominance over a strong key-aware heuristic.
